\documentclass[10pt,a4paper]{article}
\usepackage[margin=18mm]{geometry}
\usepackage{fontspec}
\usepackage{amsmath,amssymb}
\usepackage{booktabs}
\usepackage{enumitem}
\usepackage{xcolor}
\usepackage{hyperref}
\usepackage{listings}

\lstset{
  basicstyle=\ttfamily\small,
  frame=single,
  columns=fullflexible,
  breaklines=true,
  keepspaces=true
}

\setmainfont{Noto Sans CJK KR}[AutoFakeSlant=0.2]
\setsansfont{Noto Sans CJK KR}
\setmonofont{D2Coding}
\setlength{\parindent}{0pt}
\setlength{\parskip}{5pt}
\setlist{nosep,leftmargin=6mm}
\renewcommand{\arraystretch}{1.2}
\pagestyle{plain}

\NeedsTeXFormat{LaTeX2e}
\ProvidesPackage{answer_coversheet}[2026/07/30 2026 cryptanalysis contest answer cover sheet]

% Recreates the official 2026_암호분석경진대회-답안양식.hwp look: a
% "2026 암호분석경진대회" / "N번 문제" header (single rule under the first
% row, double rule under the second, matching the HWP's own
% borderfill-4/5/6 cell borders) and a page border, drawn as a background
% layer on every page. Deliberately does NOT wrap body content in any
% width-changing environment (an earlier mdframed-based version looped
% forever once a wide table met its narrowed frame width) so it never
% interacts with a document's existing tables, margins, or line widths.

\RequirePackage{geometry}
\RequirePackage{eso-pic}
\RequirePackage{tikz}
\RequirePackage{fontspec}

% Reserve room at the top of every page for the header below, and enough
% bottom clearance to stay clear of the border (19.3mm inset), without
% touching left/right margins (existing wide tables are sized against those).
\geometry{top=30mm, bottom=22mm}

\newcommand{\ProblemNumber}{}
\newcommand{\SetProblemNumber}[1]{\renewcommand{\ProblemNumber}{#1}}

% Match the font every other submission already uses for its body text
% (the HWP's real font, 함초롬바탕/Hancom HCR Batang, is a proprietary font
% not available via apt/pip here). Override with \renewcommand when the
% document's main font already has no Hangul coverage (see submissions/02).
\IfFontExistsTF{Noto Sans CJK KR}{%
  \newfontfamily\AnswerCoverFont{Noto Sans CJK KR}[AutoFakeSlant=0.2]%
}{%
  \newcommand{\AnswerCoverFont}{}%
}

% Box/row measurements below are taken directly from the HWP's own layout
% (via hwp5html): HeaderPageFooter margin-top 10mm + page padding-top 5mm =
% 15mm top inset; margin-bottom 5mm + padding-bottom 5mm = 10mm, plus the
% content cell's own extra spacing brings it to ~19.3mm; table margin-left
% 1mm inside a 10mm side margin = 11mm; each header row is 6.15mm tall.
% The title/problem-number header (and its rules) is drawn on page 1 only;
% the outer border repeats on every page.
\AddToShipoutPictureBG{%
  \AnswerCoverFont
  \begin{tikzpicture}[overlay]
    \coordinate (TL) at ([xshift=11mm]current page.north west);
    \coordinate (TR) at ([xshift=-11mm]current page.north east);
    \coordinate (BR) at ([xshift=-11mm]current page.south east);
    \draw[line width=0.6pt] ([yshift=-15mm]TL) rectangle ([yshift=19.3mm]BR);
    \ifnum\value{page}=1
      \draw[line width=0.6pt] ([yshift=-21.15mm]TL) -- ([yshift=-21.15mm]TR);
      \draw[line width=0.4pt] ([yshift=-27.1mm]TL) -- ([yshift=-27.1mm]TR);
      \draw[line width=0.4pt] ([yshift=-27.5mm]TL) -- ([yshift=-27.5mm]TR);
      \node[anchor=north] at ([yshift=-15.5mm]current page.north) {2026 암호분석경진대회};
      \node[anchor=north] at ([yshift=-21.65mm]current page.north) {\ProblemNumber 번 문제};
    \fi
  \end{tikzpicture}%
}

\endinput

\SetProblemNumber{5}

\title{동형암호}
\author{비밀 다항식 및 고정 보고문 복원}
\date{2026년 7월}

\begin{document}
\maketitle

\section*{결과 요약}

복구된 고정 보고문(State)의 평문은 다음과 같다.
\begin{center}
\fbox{\texttt{BGV DAILY STATUS CORE-A LINK OK TEMP NORMAL POWER STABLE}}
\end{center}

비밀 다항식 $s$ 및 패딩을 제외한 상태 보고문(State)의 계수는 각각 \texttt{02\_secret\_s.txt}, \texttt{03\_state.txt}에 $x^0$부터 차수 \\
오름차순으로 기록하였으며, 최종 복원된 날짜는 \texttt{20260410} 및 \texttt{20260411}이고 실제 NUL 패딩의 길이는 0바이트이다.

\section{공약수 취약점과 노이즈 제거}

textbook-BGV 암호문은 환 $R_q=\mathbb Z_q[x]/(x^{64}+1)$ 상에서 다음 식을 만족한다.
\[
  c_1=c_0s+te+m\pmod q
\]
두 연속된 날짜의 암호문을 감산(차분)하면 다음 관계식을 얻는다.
\[
  \Delta c_1=\Delta c_0s+t\Delta e+\Delta m\pmod q
\]
주어진 모듈러스 매개변수에는 다음과 같은 결정적인 대수적 결함이 존재한다.
\[
  \boxed{\gcd(q,t)=257}
\]
공약수인 257은 $q$와 $t$를 모두 나누므로, 양변을 유한체 $\mathbb F_{257}$ 상으로 축소(reduction)하면 오류 항 $t\Delta e$ 및 modulo-$q$ 관계식이 소거되어 노이즈가 완벽히 배제된다.
\[
  \boxed{\Delta c_1=\Delta c_0s+\Delta m\pmod {257}}
\]

\section{날짜 차분 전수조사와 삼원 비밀 다항식 (Ternary Secret) 복원}

두 날짜 평문의 상태 보고문(State) 및 중간 NUL 패딩 구조가 동일하므로, 메시지 차분 다항식 $\Delta m$의 앞 56개 계수는 0이며 마지막 8개 계수만이 날짜 문자열(\texttt{YYYYMMDD})의 일자 변동에 따른 차분값이다.

대상 연도의 범위를 제한하지 않기 위해 네 자리 그레고리력 연도 \texttt{0001..9999}에서 발생 가능한 모든 날짜 변동 패턴을 탐색 대상에 포함하였다. 날짜 전이는 동일 월 내에서의 전환(\texttt{01} $\to$ \texttt{02} 등)과 월말에서 익월 1일로 넘어가는 전환으로 분류된다. 이를 모두 생성한 뒤 중복을 제거하면, modulo-257 상에서 서로 다른 날짜 차분 패턴은 최종적으로 11개로 축소된다.

각 패턴에 대해 $\Delta c_0$와의 \emph{negacyclic} convolution 행렬 $A\in\mathbb F_{257}^{64\times64}$를 정의하고, $As\equiv\Delta c_1-\Delta m\pmod{257}$ 선형계를 구성한 뒤 가우스 소거법(Gaussian Elimination)을 통해 해를 구하였다. 고유 해의 유일성을 입증하기 위해, \\
최초의 삼원(ternary) 해 발견 시 탐색을 중단하지 않고 전체 affine 해공간(257개 후보)을 전수 조사하였다.

\begin{center}
\begin{tabular}{lr}
\toprule
검증량 & 결과 \\
\midrule
서로 다른 다음 날 delta 패턴 수 & 11 \\
정답 선형계 rank / nullity & 63 / 1 \\
정답 affine space 전수조사 수 & 257개 \\
모든 계수가 $\{-1,0,1\}$인 대수 후보 수 & 1개 \\
날짜·State·noise 검증 통과 수 & 1개 \\
\bottomrule
\end{tabular}
\end{center}

\section{후보 다항식의 유일성 및 원본 암호문 검증}

수집한 모든 ternary 후보를 두 원본 ciphertext에 역적용하여 다음 대수적 및 논리적 조건을 동시에 검증하였다.

\begin{itemize}
  \item 복원된 날짜가 8자리 ASCII 형식의 유효한 그레고리력 날짜이며, 두 날짜가 정확히 하루 차이이다.
  \item 두 \texttt{State || 00...00} prefix가 동일하고, 우측 NUL 패딩을 제거한 상태 보고문 문자열이 출력 가능한 ASCII 문자로 구성된다.
  \item 두 날 모두 $c_1=c_0s+te+m\pmod q$를 만족하며, 복원된 에러 다항식 $e$의 계수 분포가 BGV 파라미터 한계 내인 $1 \sim 4$ 범위에 존재한다.
\end{itemize}

최종 복원된 날짜는 \texttt{20260410} 및 \texttt{20260411}이며, 실제 NUL 패딩 길이는 0바이트이다. 비밀 다항식 $s$ 및 패딩을 제외한 상태 보고문(State)의 계수는 각각 \texttt{02\_secret\_s.txt}, \texttt{03\_state.txt}에 $x^0$부터 차수 오름차순으로 기록하였다.

\section{재현 및 코드 검증}

제출 패키지에 포함된 검증용 스크립트인 \texttt{solve\_bgv.py}는 다음 명령어로 실행 가능하며, 실행 시 두 제출 \\
 다항식 파일(\texttt{02\_secret\_s.txt}, \texttt{03\_state.txt})과의 일치 여부 및 원 암호문의 무결성을 검증하고 자가 진단(Self-Check PASS)을 수행한다.

\begin{lstlisting}
python3 solve_bgv.py
\end{lstlisting}

\section{참고 자료}

\begin{itemize}
  \item Brakerski et al.,
    \href{https://eprint.iacr.org/2011/277}
    {\emph{Fully Homomorphic Encryption without Bootstrapping}}: Leveled FHE 및 RLWE 기반 BGV 동형암호의 노이즈 관리 기법 연구.
  \item HomomorphicEncryption.org,
    \href{https://homomorphicencryption.org/standard/}
    {\emph{Homomorphic Encryption Standard}}: 평문/암호문 모듈러스, 다항식 환 구조 및 에러 표기법 표준 정의.
  \item Lyubashevsky et al.,
    \href{https://eprint.iacr.org/2012/230}
    {\emph{On Ideal Lattices and Learning with Errors over Rings}}: Ring-LWE 문제의 오류 분포 및 이상 격자(Ideal Lattice) 구조 분석.
\end{itemize}

\end{document}
